% =================================================================
% Sprint 1 - Manage Forum
% The functions below follow the detailed form of the report template:
% function trigger, function description, screen layout and function details.
% They are written against the implementation in
%   apps/api/src/routes/forums/forums-post, forums-comment, forums-like,
%   forums-reports and apps/web/app/forum, including the BFF proxy
%   (apps/web/app/api/backend) that sends post creation through the
%   SensitiveAI moderation proxy.
% =================================================================

\subsection{Manage Forum}

\subsubsection{Function: View Forum Category List}

\fnhead{Function Trigger}
Any visitor opens the Community Forum area from the main navigation, or follows a direct
link to the forum landing page. The function is public, therefore it runs on every visit of
the page, whether or not the visitor is signed in.\newline
\textbf{Navigation path:} Home $\rightarrow$ Forum (\texttt{/forum}).

\fnhead{Function Description}

\sublabel{Actors / Roles}
\begin{itemize}[leftmargin=1.6em,topsep=3pt,itemsep=2pt,parsep=0pt,label=$\bullet$]
    \item \textbf{Primary Actor:} Guest, Freelancer, Client and Administrator. The screen is
          readable by every role, including visitors without an account.
    \item \textbf{Supporting Systems:} Frevia Web application (Next.js), Frevia BFF proxy
          (\texttt{/api/backend}), Frevia REST API (NestJS) and PostgreSQL accessed through
          Prisma ORM.
\end{itemize}

\sublabel{Purpose}
\noindent The function is the entry point of the community area. It presents the active forum
categories together with the number of published posts that each one contains, and it
highlights both the categories and the members with the highest activity, so a visitor can
reach an ongoing discussion in one click and can immediately see which topics are alive.

\sublabel{Interface}
\begin{itemize}[leftmargin=1.6em,topsep=3pt,itemsep=2pt,parsep=0pt,label=$\bullet$]
    \item Community Forum page at \texttt{/forum}, opened from the main navigation.
    \item Grid of category cards; each card shows the category name, its description and the
          number of approved posts.
    \item ``Most Active'' panel listing the three categories with the largest number of posts.
    \item ``Top Contributors'' panel listing the five members with the largest number of posts,
          with their post and comment counters.
    \item Empty-state panel displayed when no category has been created yet.
\end{itemize}

\sublabel{Data Processing Overview}
\begin{itemize}[leftmargin=1.6em,topsep=3pt,itemsep=2pt,parsep=0pt,label=$\bullet$]
    \item The page requests \texttt{GET /api/forums/categories}.
    \item The API keeps only the categories whose \texttt{deletedAt} is null and orders the
          result by \texttt{name} in ascending order.
    \item For every category the API counts the posts that are not soft deleted and whose
          \texttt{moderationStatus} is \texttt{APPROVED}; the count is returned as
          \texttt{postCount}, so posts awaiting moderation or rejected by it are never counted.
    \item The ``Most Active'' panel requests \texttt{GET /api/forums/categories/top?limit=3};
          the API sorts the same projection by post count in descending order.
    \item The ``Top Contributors'' panel requests \texttt{GET /api/forums/users/top?limit=5};
          the API selects the users that are neither deleted nor banned, counts their posts and
          comments with \texttt{deletedAt} null and orders them by post count in descending order.
    \item The category list is cached by the client for thirty minutes and is re-used while the
          visitor moves between forum screens, so returning to the list does not repeat the call.
    \item A category without a description is rendered with the fallback text ``No description
          provided.''; an empty result set is rendered as ``No categories yet''.
\end{itemize}

\fnhead{Screen Layout (Mock-up Description)}
\screenfigure[0.95]{image/screenshot/sprint1/kvy/ViewCategoriesList.png}{Forum category list screen}
\begin{itemize}[leftmargin=1.6em,topsep=3pt,itemsep=2pt,parsep=0pt,label=$\bullet$]
    \item The page header states the title ``Community Forum'' and a short description of the
          area, which introduces the forum to first-time visitors.
    \item Each category card is a link to the category detail screen and displays the category
          name, the description and the number of published posts.
    \item The ``Most Active'' panel ranks categories by post volume and gives a shortcut to the
          busiest discussions.
    \item The ``Top Contributors'' panel shows the avatar, the display name, the post count and
          the comment count of the most active members.
    \item When the list is empty, a placeholder explains that categories will appear once they
          are created, which happens after an administrator creates them.
\end{itemize}

\fnhead{Function Details}

\sublabel{Input Data}
\begin{specdatatable}{Data Name}{Description}
Page route & \texttt{/forum}. No parameter is required to display the list.\\ \hline
\texttt{limit} & Number of records returned by the two ranking panels; integer between 1 and
10, default 3. It is only used by the ``Most Active'' and ``Top Contributors'' calls.\\ \hline
Visitor session & Optional access token forwarded by the BFF proxy. A guest receives the same
list; the ranking panels do not expose any personal action.\\ \hline
\end{specdatatable}

\sublabel{Validation Rules}
\begin{itemize}[leftmargin=1.6em,topsep=3pt,itemsep=2pt,parsep=0pt,label=$\bullet$]
    \item The list endpoint takes no parameter, therefore it cannot be malformed; it always
          returns the complete list of visible categories.
    \item The \texttt{limit} of both ranking endpoints is coerced to an integer, must be at
          least 1 and at most 10, and defaults to 3 when it is absent.
    \item A category is visible only when its \texttt{deletedAt} column is null; a category moved
          to the trash is excluded from the list and from both rankings.
    \item A post contributes to a counter only when it is not soft deleted and its
          \texttt{moderationStatus} is \texttt{APPROVED}.
    \item A user appears in ``Top Contributors'' only when the account is not deleted and is not
          banned.
\end{itemize}

\sublabel{Business Rules}
\begin{itemize}[leftmargin=1.6em,topsep=3pt,itemsep=2pt,parsep=0pt,label=$\bullet$]\brrulelist
    \item BR-FORUM-01: Only categories with \texttt{deletedAt = NULL} shall be displayed.
    \item BR-FORUM-02: Categories shall be ordered by \texttt{name} in ascending order.
    \item BR-FORUM-03: \texttt{postCount} shall count only posts with \texttt{deletedAt = NULL}
          and \texttt{moderationStatus = APPROVED}.
\end{itemize}

\sublabel{Normal Flow}
\begin{enumerate}[leftmargin=1.8em,topsep=3pt,itemsep=2pt,parsep=0pt]
    \item The visitor opens the Forum area from the navigation menu.
    \item The web application renders the page shell and issues the category request.
    \item The BFF proxy forwards the request to the REST API, adding the access token when the
          visitor is signed in.
    \item The API queries the categories that are not soft deleted and counts their approved
          posts.
    \item The API returns the ordered list of categories with \texttt{postCount}.
    \item The client renders the category grid, applying the description fallback where needed.
    \item The client issues the two ranking requests for the side panels.
    \item The API returns the three most active categories and the five most active members.
    \item The client renders both panels; the page is now fully displayed.
    \item The visitor selects a category card and is routed to the category detail screen.
\end{enumerate}

\sublabel{Abnormal Flow / Exception Handling}
\begin{flowtable}{Case}{Description}
A1 & The database cannot be queried. The API returns a server error carrying
\texttt{Error.FailedToLoadForumCategories} and the page shows a generic error instead of the
grid.\\ \hline
A2 & No category exists or every category has been soft deleted. The page renders the empty
state ``No categories yet''; the ranking panels are hidden.\\ \hline
A3 & A category has no description. The card shows ``No description provided.'' instead of an
empty area.\\ \hline
A4 & The ranking call fails while the list succeeds. The grid remains usable and only the
affected panel is omitted, so navigation is never blocked.\\ \hline
A5 & The visitor loses the network connection. The client displays the generic error message
and offers a retry by reloading the page.\\ \hline
A6 & The access token has expired. The BFF proxy refreshes the token once and repeats the call;
if the refresh fails the visitor continues as a guest.\\ \hline
A7 & The account was banned while the page was open. The BFF proxy detects the ban response,
clears the authentication cookies and the client redirects to the sign-in page.\\ \hline
\end{flowtable}

\sublabel{Output Data}
\begin{specdatatable}{Data Name}{Description}
\texttt{categories[]} & Visible categories ordered by name, each with \texttt{id},
\texttt{name}, \texttt{slug}, \texttt{description}, \texttt{createdAt}, \texttt{updatedAt} and
\texttt{postCount}.\\ \hline
\texttt{categories} (top) & Up to three categories ordered by post count in descending order,
with the same fields as the list.\\ \hline
\texttt{users} (top) & Up to five members with \texttt{id}, \texttt{displayName},
\texttt{avatarUrl}, \texttt{postCount} and \texttt{commentCount}.\\ \hline
Empty collection & An empty array when no category is visible; the client renders the empty
state described above.\\ \hline
\end{specdatatable}

\sublabel{Post-Conditions}
\begin{itemize}[leftmargin=1.6em,topsep=3pt,itemsep=2pt,parsep=0pt,label=$\bullet$]
    \item The visitor sees every active category with the number of published posts it holds.
    \item No data is modified: the function is read-only.
    \item Deleted categories, deleted posts, rejected posts and pending posts are never exposed
          through this screen.
\end{itemize}

% =================================================================

\subsubsection{Function: View Forum Category Detail and Post List}

\fnhead{Function Trigger}
The visitor selects a category card on the forum landing page, follows a shared category link,
or searches for a post from the category screen.\newline
\textbf{Navigation path:} Home $\rightarrow$ Forum $\rightarrow$ \textit{Category}
(\texttt{/forum/\{category-slug\}-\{category-id\}}).

\fnhead{Function Description}

\sublabel{Actors / Roles}
\begin{itemize}[leftmargin=1.6em,topsep=3pt,itemsep=2pt,parsep=0pt,label=$\bullet$]
    \item \textbf{Primary Actor:} Guest, Freelancer, Client and Administrator.
    \item \textbf{Supporting Systems:} Frevia Web application, Frevia BFF proxy, Frevia REST API
          and PostgreSQL through Prisma ORM.
\end{itemize}

\sublabel{Purpose}
\noindent The function lists the published posts of one category so that members can read,
search and filter the discussions of a topic. It also gives signed-in members the entry point
to create a post, to restrict the list to their own posts, and to reach the most discussed
posts of the week.

\sublabel{Interface}
\begin{itemize}[leftmargin=1.6em,topsep=3pt,itemsep=2pt,parsep=0pt,label=$\bullet$]
    \item Breadcrumb Home $\rightarrow$ Forum $\rightarrow$ \textit{Category} and the category
          header with its name, description and post count.
    \item Search box with the placeholder ``Search posts\ldots'' and a clear-search button.
    \item ``My Posts'' toggle, displayed only to signed-in members.
    \item ``New Post'' button, which opens the post creation dialog; a guest is routed to the
          sign-in page instead.
    \item List of post cards with the title, an excerpt of the content, the author, the
          creation date, the like count and the comment count.
    \item Page-size selector, page numbers with ellipsis, and the ``Showing x of y'' summary.
    \item ``Trending in \{category\}'' panel listing the three most discussed posts of the last
          seven days.
    \item Empty states for no post, no search result and no personal post.
\end{itemize}

\sublabel{Data Processing Overview}
\begin{itemize}[leftmargin=1.6em,topsep=3pt,itemsep=2pt,parsep=0pt,label=$\bullet$]
    \item The page resolves the category from the slug in the URL and requests
          \texttt{GET /api/forums/posts?page=\ldots\&limit=\ldots\&categoryId=\ldots}.
    \item When the visitor submits the search box, the client navigates to the same route with a
          \texttt{search} parameter and resets the page to 1.
    \item When ``My Posts'' is active, the client adds the identifier of the signed-in member as
          the \texttt{userId} parameter.
    \item The API filters the posts whose \texttt{deletedAt} is null and whose
          \texttt{moderationStatus} is \texttt{APPROVED}, applies the category, author and title
          filters, and sorts the result by \texttt{createdAt} in descending order.
    \item Lists and counts are read in one database transaction, then the API returns the page
          together with the pagination block.
    \item The trending panel requests the top posts of the category for the last seven days,
          ordered by like count and then by comment count.
    \item Changing the page, the page size, the search term or the ``My Posts'' filter rewrites
          the query string, so the state of the list can be shared and restored by reloading.
\end{itemize}

\fnhead{Screen Layout (Mock-up Description)}
\screenfigure[0.95]{image/screenshot/sprint1/kvy/ViewCategoryDetail+ViewPostList.png}{Category detail screen with the list of its posts}
\begin{itemize}[leftmargin=1.6em,topsep=3pt,itemsep=2pt,parsep=0pt,label=$\bullet$]
    \item The breadcrumb keeps the visitor oriented between the forum landing page, the current
          category and the post that will be opened.
    \item The category header repeats the name and the description, which are maintained by the
          administrator who created the category.
    \item The toolbar groups the search box, the ``My Posts'' toggle and the ``New Post''
          button; the search box submits on Enter and can be cleared with the cross icon.
    \item Every card shows the post title, the beginning of its content without the rich-text
          markup, the author avatar and name, the date, and the like and comment counters.
    \item Own posts are tagged with ``(yours)'' so a member can recognise them in the list.
    \item The footer of the list shows the pagination and the range of records currently
          displayed.
\end{itemize}

\fnhead{Function Details}

\sublabel{Input Data}
\begin{specdatatable}{Data Name}{Description}
\texttt{categoryId} & Identifier of the category taken from the URL; it restricts the list to
one category.\\ \hline
\texttt{page} & Page number, integer greater than or equal to 1, default 1.\\ \hline
\texttt{limit} & Page size, integer greater than or equal to 1, default 10; the client offers a
selector and re-sends the request when it changes.\\ \hline
\texttt{search} & Optional search term applied to the post title, case insensitive.\\ \hline
\texttt{userId} & Optional author filter sent when ``My Posts'' is active; it always carries the
identifier of the signed-in member.\\ \hline
\end{specdatatable}

\sublabel{Validation Rules}
\begin{itemize}[leftmargin=1.6em,topsep=3pt,itemsep=2pt,parsep=0pt,label=$\bullet$]
    \item \texttt{page} and \texttt{limit} are coerced to integers and must be at least 1;
          the default values are applied by the API when the parameters are absent.
    \item \texttt{categoryId} and \texttt{userId} must be positive integers when they are sent.
    \item \texttt{search} is optional and is compared with the title only, without case
          sensitivity; leading and trailing spaces are trimmed by the client.
    \item A requested page beyond the last page returns an empty list while the pagination block
          still reports the real total, so the client can render an empty state.
    \item Only approved and non-deleted posts are counted and displayed, therefore the totals
          never include posts that are pending, rejected or deleted.
\end{itemize}

\sublabel{Business Rules}
\begin{itemize}[leftmargin=1.6em,topsep=3pt,itemsep=2pt,parsep=0pt,label=$\bullet$]\brrulelist
    \item BR-FORUM-06: Only posts with \texttt{deletedAt = NULL} and
          \texttt{moderationStatus = APPROVED} shall be displayed.
    \item BR-FORUM-07: Posts shall be sorted by \texttt{createdAt} in descending order.
    \item BR-FORUM-08: The system shall support pagination and configurable page sizes.
    \item BR-FORUM-09: \texttt{commentCount} shall count only comments with
          \texttt{deletedAt = NULL}.
\end{itemize}

\sublabel{Normal Flow}
\begin{enumerate}[leftmargin=1.8em,topsep=3pt,itemsep=2pt,parsep=0pt]
    \item The visitor selects a category card on the forum landing page.
    \item The client routes to the category screen and reads the page, the page size and the
          optional filters from the query string.
    \item The client requests the posts of the category from the API.
    \item The API validates the filter, resolves the pagination offset and runs the list and
          count queries in one transaction.
    \item The API applies the visibility rules and orders the posts by creation date from the
          newest to the oldest.
    \item The API returns the page of posts with their like and comment counters plus the
          pagination block.
    \item The client renders the cards, strips the rich-text markup for the excerpt and applies
          the empty states when no card is returned.
    \item The client requests the trending posts of the category and renders the side panel.
    \item The visitor optionally searches, changes the page size, moves to another page or
          activates ``My Posts''; each action re-issues the request with the new parameters.
    \item The visitor selects a post card and is routed to the post detail screen.
\end{enumerate}

\sublabel{Abnormal Flow / Exception Handling}
\begin{flowtable}{Case}{Description}
A1 & The category does not exist or was soft deleted. The API returns
\texttt{Error.ForumCategoryNotFound} and the client displays the not-found state with a link
back to the forum.\\ \hline
A2 & The database query fails. The API returns \texttt{Error.FailedToLoadForumPosts} and the
list area shows a generic error instead of the cards.\\ \hline
A3 & The requested page is larger than the last page. The API returns an empty list and the
client renders ``No posts found'' with the pagination intact.\\ \hline
A4 & The search term matches no approved post. The client reports ``No results for
``\{query\}''. Try a different search.'' and keeps the clear-search action available.\\ \hline
A5 & ``My Posts'' is activated by a guest. The toggle is not rendered; if the parameter is
forced in the URL the API answers with its validation error and the list is not displayed.\\ \hline
A6 & The page size or page number is not a positive integer. The validation pipe rejects the
request before the query is executed and the client falls back to the first page.\\ \hline
A7 & The trending request fails. The list remains fully usable and only the side panel is
omitted.\\ \hline
A8 & The visitor loses the connection. The client shows the generic error message and the
visitor can retry by reloading the current URL.\\ \hline
\end{flowtable}

\sublabel{Output Data}
\begin{specdatatable}{Data Name}{Description}
\texttt{posts[]} & Visible posts of the requested page with \texttt{id}, \texttt{categoryId},
\texttt{userId}, \texttt{title}, \texttt{slug}, \texttt{content}, \texttt{moderationStatus},
\texttt{createdAt}, \texttt{updatedAt}, \texttt{likeCount}, \texttt{commentCount} and the author
identity.\\ \hline
\texttt{pagination} & \texttt{page}, \texttt{limit}, \texttt{total} and \texttt{totalPages}
computed from the filtered count.\\ \hline
\texttt{posts[]} (top) & Trending posts of the category with the interaction score used for the
ranking.\\ \hline
Empty collection & An empty array when the page, the filter or the search matches no approved
post.\\ \hline
\end{specdatatable}

\sublabel{Post-Conditions}
\begin{itemize}[leftmargin=1.6em,topsep=3pt,itemsep=2pt,parsep=0pt,label=$\bullet$]
    \item The visitor sees one page of approved posts of the selected category, newest first.
    \item The search, the page size and the author filter are reflected in the URL and can be
          restored by reloading the page.
    \item No data is modified: the function is read-only.
\end{itemize}

% =================================================================

\subsubsection{Function: View Post Detail}

\fnhead{Function Trigger}
The visitor selects a post card, follows a shared post link, or opens a related post from the
side panel of another discussion.\newline
\textbf{Navigation path:} Home $\rightarrow$ Forum $\rightarrow$ \textit{Category} $\rightarrow$
\textit{Post} (\texttt{/forum/\{category\}/\{post-slug\}-\{post-id\}}).

\fnhead{Function Description}

\sublabel{Actors / Roles}
\begin{itemize}[leftmargin=1.6em,topsep=3pt,itemsep=2pt,parsep=0pt,label=$\bullet$]
    \item \textbf{Primary Actor:} Guest, Freelancer, Client and Administrator.
    \item \textbf{Supporting Systems:} Frevia Web application, Frevia BFF proxy, Frevia REST API
          and PostgreSQL through Prisma ORM.
\end{itemize}

\sublabel{Purpose}
\noindent The function displays a complete discussion: the post with its formatted content and
its author, the engagement counters, the list of members who liked it and the comment thread.
It is the screen from which every write action on the discussion is started, and it shows the
author-specific actions only to the author.

\sublabel{Interface}
\begin{itemize}[leftmargin=1.6em,topsep=3pt,itemsep=2pt,parsep=0pt,label=$\bullet$]
    \item Breadcrumb Home $\rightarrow$ Forum $\rightarrow$ \textit{Category} $\rightarrow$
          \textit{Post title} and the line ``Discussion in \textit{Category}''.
    \item Post title, author avatar, display name, creation date and the ``Edited \{date\}''
          marker when the post was modified after its creation.
    \item Rich-text body rendered with headings, lists and links.
    \item Action bar with the like button and its counter, the comment counter, the ``Report''
          action and, for the author, the ``Edit'' and ``Delete'' actions.
    \item Likes dialog and comment section described by the next functions.
    \item ``Related Posts'' side panel with other discussions of the same category.
\end{itemize}

\sublabel{Data Processing Overview}
\begin{itemize}[leftmargin=1.6em,topsep=3pt,itemsep=2pt,parsep=0pt,label=$\bullet$]
    \item The client resolves the post from the slug and requests
          \texttt{GET /api/forums/posts/by-slug/\{slug\}}.
    \item The API resolves the identifier of the post and returns the detail projection with the
          category summary and the author identity.
    \item Only a post that is not soft deleted and whose status is \texttt{APPROVED} can be
          resolved; a pending or rejected post answers with a not-found error.
    \item The client requests the likes of the post to compute the counter and to derive whether
          the signed-in member already liked it.
    \item The client requests the first fifty comments of the post to render the thread and the
          comment counter.
    \item The client requests the related posts of the category for the side panel.
    \item The page compares \texttt{updatedAt} with \texttt{createdAt} to decide whether the
          ``Edited'' marker is displayed, and compares the author identifier with the session
          identifier to decide whether the edit and delete actions are rendered.
\end{itemize}

\fnhead{Screen Layout (Mock-up Description)}
\screenfigure[0.95]{image/screenshot/sprint1/kvy/ViewPostDetail.png}{Post detail screen with the comment section}
\begin{itemize}[leftmargin=1.6em,topsep=3pt,itemsep=2pt,parsep=0pt,label=$\bullet$]
    \item The title block gives the discussion context and the path back to its category.
    \item The author block shows the avatar, the display name and the publication date; when the
          post was edited, the edition date is added next to the creation date.
    \item The body is rendered as formatted content, so headings, bullet lists and links written
          in the editor keep their appearance.
    \item The action bar holds the heart button with the like counter, the comment counter and
          the ``Report'' action; the heart is filled once the member liked the post.
    \item The edit and delete actions appear only for the author of the post.
    \item The comment section below the action bar lists the thread and gives the signed-in
          member the comment editor.
    \item The right side panel proposes related posts of the same category.
\end{itemize}

\fnhead{Function Details}

\sublabel{Input Data}
\begin{specdatatable}{Data Name}{Description}
\texttt{slug} & Post slug taken from the URL; it uniquely identifies the discussion.\\ \hline
\texttt{postId} & Identifier resolved from the slug and used by the likes, comments and report
calls of the same screen.\\ \hline
\texttt{page}, \texttt{limit} & Pagination of the comment thread; the detail screen requests the
first page with a size of 50.\\ \hline
Session identity & Identifier of the signed-in member, used to compute the liked state and to
decide whether the author actions are displayed.\\ \hline
\end{specdatatable}

\sublabel{Validation Rules}
\begin{itemize}[leftmargin=1.6em,topsep=3pt,itemsep=2pt,parsep=0pt,label=$\bullet$]
    \item The slug must resolve to a post that is not soft deleted and whose
          \texttt{moderationStatus} is \texttt{APPROVED}; otherwise the API answers with
          \texttt{Error.ForumPostNotFound}.
    \item The detail projection must contain a title and a content, which are required fields of
          the post model.
    \item The comment pagination follows the same rules as the post list: integers greater than
          or equal to 1, default page 1 and default size 10.
    \item The liked state is derived only from the likes that belong to the post and match the
          session identifier; a guest always sees the unliked state.
    \item The author actions are rendered only when the session identifier equals the post
          author identifier, never from a value supplied by the client.
\end{itemize}

\sublabel{Business Rules}
\begin{itemize}[leftmargin=1.6em,topsep=3pt,itemsep=2pt,parsep=0pt,label=$\bullet$]\brrulelist
    \item BR-FORUM-10: Only posts with \texttt{deletedAt = NULL} and
          \texttt{moderationStatus = APPROVED} shall be displayed.
    \item BR-FORUM-11: Post details shall include title, content, author, category and created
          date.
    \item BR-FORUM-09: \texttt{commentCount} shall count only comments with
          \texttt{deletedAt = NULL}.
    \item BR-FORUM-28: The like detail shall include the user identifier and the profile of each
          member who liked the post.
\end{itemize}

\sublabel{Normal Flow}
\begin{enumerate}[leftmargin=1.8em,topsep=3pt,itemsep=2pt,parsep=0pt]
    \item The visitor selects a post card in the category screen.
    \item The client extracts the slug from the URL and requests the post detail.
    \item The API resolves the slug to a post that is visible and returns the detail projection.
    \item The client renders the breadcrumb, the title, the author block and the formatted body.
    \item The client requests the likes of the post and the first page of comments.
    \item The API returns both collections with the profile of their authors.
    \item The client computes the like counter, the liked state and the comment counter.
    \item The client renders the action bar and exposes the author actions only to the author.
    \item The client requests the related posts of the category and renders the side panel.
    \item The visitor reads the discussion and can like the post, open the likes list, write a
          comment, report the content or, as the author, edit or delete the post.
\end{enumerate}

\sublabel{Abnormal Flow / Exception Handling}
\begin{flowtable}{Case}{Description}
A1 & The slug does not resolve to any post. The client displays the not-found screen with the
message ``This post may have been deleted or the link is incorrect.'' and the action ``Back to
category''.\\ \hline
A2 & The post is still pending moderation or was rejected. The API hides it, therefore the same
not-found screen is displayed and the post is never exposed before approval.\\ \hline
A3 & The post was soft deleted by its author or by an administrator. The detail call answers with
\texttt{Error.ForumPostNotFound} and the not-found screen is displayed.\\ \hline
A4 & The likes request fails. The post body remains readable, the like counter falls back to
zero and the action is disabled.\\ \hline
A5 & The comments request fails. The thread area shows an error message while the post body and
the action bar remain usable.\\ \hline
A6 & The session expires while the page is open. The BFF proxy refreshes the token once; when
the refresh fails the visitor continues as a guest without author actions.\\ \hline
A7 & The visitor follows a link of a deleted category. The category is not found, so the client
shows the not-found screen instead of a broken layout.\\ \hline
A8 & The connection drops while the page loads. The client shows the generic error message and
the visitor can retry by reloading.\\ \hline
\end{flowtable}

\sublabel{Output Data}
\begin{specdatatable}{Data Name}{Description}
Post detail & \texttt{id}, \texttt{categoryId}, \texttt{userId}, \texttt{title},
\texttt{slug}, \texttt{content}, \texttt{moderationStatus}, \texttt{moderationScore},
\texttt{createdAt} and \texttt{updatedAt}.\\ \hline
\texttt{category} & Identifier, name and slug of the category the post belongs to, or null when
the post has no category.\\ \hline
\texttt{user} & Identifier and profile of the author, with the display name and the avatar.\\ \hline
\texttt{likes[]} & Members who liked the post, with their identifier, reaction type, creation date
and profile.\\ \hline
\texttt{comments[]} & First page of the thread with the content, the author, the dates, the like
counter and the liked state of the current member.\\ \hline
\texttt{pagination} & Total number of visible comments and the pagination of the requested page.\\ \hline
\end{specdatatable}

\sublabel{Post-Conditions}
\begin{itemize}[leftmargin=1.6em,topsep=3pt,itemsep=2pt,parsep=0pt,label=$\bullet$]
    \item The visitor reads the complete discussion with its engagement counters.
    \item The action bar exposes only the actions allowed by the role and the ownership of the
          post.
    \item No data is modified by the read itself; the counters reflect only visible rows.
\end{itemize}

% =================================================================

\subsubsection{Function: Create Post}

\fnhead{Function Trigger}
A signed-in member selects the ``New Post'' button in the toolbar of a category screen, writes
the post and confirms with the ``Post'' action. The button is rendered for Clients, Freelancers
and Administrators; a guest who follows the button is routed to the sign-in page instead.\newline
\textbf{Navigation path:} Home $\rightarrow$ Forum $\rightarrow$ \textit{Category} $\rightarrow$
New Post.

\fnhead{Function Description}

\sublabel{Actors / Roles}
\begin{itemize}[leftmargin=1.6em,topsep=3pt,itemsep=2pt,parsep=0pt,label=$\bullet$]
    \item \textbf{Primary Actor:} Freelancer and Client. An Administrator uses the same dialog
          when browsing the forum.
    \item \textbf{Supporting Systems:} Frevia Web application, Frevia BFF proxy, SensitiveAI
          moderation proxy, Frevia REST API and PostgreSQL through Prisma ORM.
\end{itemize}

\sublabel{Purpose}
\noindent The function lets a member start a new discussion inside a category. Because the
content is user generated, every submission is screened by the moderation proxy before it
reaches the forum, and the resulting post is stored with the moderation verdict so that only
approved discussions become visible to the community.

\sublabel{Interface}
\begin{itemize}[leftmargin=1.6em,topsep=3pt,itemsep=2pt,parsep=0pt,label=$\bullet$]
    \item Modal dialog titled ``Create Post in \textit{Category}'' with the description
          ``Share your thoughts with the community.''
    \item ``Title'' input with the placeholder ``What's on your mind?''.
    \item ``Content'' rich-text editor with the placeholder ``Write your post content here\ldots''
          and the usual formatting tools.
    \item ``Cancel'' action, which discards the draft, and ``Post'' action, which submits it and
          shows a spinner while the request is running.
    \item Toast notification after the submission, either the confirmation of publication or the
          message that the post awaits moderator review.
\end{itemize}

\sublabel{Data Processing Overview}
\begin{itemize}[leftmargin=1.6em,topsep=3pt,itemsep=2pt,parsep=0pt,label=$\bullet$]
    \item The dialog trims the title and the content and blocks the submission while either of
          them is empty.
    \item The client posts the form to \texttt{/api/backend/forums/posts}.
    \item The BFF proxy recognises this exact route as the only forum route that must be screened
          and forwards it to the SensitiveAI moderation proxy instead of the REST API.
    \item The moderation proxy scans the text, attaches the verdict (status, score and categories)
          to the request body and forwards the request to the REST API.
    \item The API validates the body, verifies that the category exists and is not soft deleted,
          and creates the post.
    \item The slug is generated from the title, normalised to lower case without accents, and
          completed with a timestamp and a random suffix to keep it unique.
    \item The moderation status is taken from the verdict; when the proxy returns no verdict the
          post is stored as \texttt{PENDING}, which means it is invisible until an administrator
          approves it.
    \item The API returns the created post and the client invalidates the post lists of the
          category so the next visit shows the new discussion.
    \item When the returned status is \texttt{PENDING}, the client shows the moderation message and
          does not navigate, because the post detail would not be reachable yet; otherwise it
          navigates to the new post.
\end{itemize}

\fnhead{Screen Layout (Mock-up Description)}
\screenfigure[0.85]{image/screenshot/sprint1/kvy/CreatePost.png}{Create post dialog}
\begin{itemize}[leftmargin=1.6em,topsep=3pt,itemsep=2pt,parsep=0pt,label=$\bullet$]
    \item The dialog title names the category in which the post will be created, so the member
          cannot post in the wrong place by mistake.
    \item The title field is a single-line input; the content field is a rich-text editor that
          stores formatted HTML.
    \item The ``Post'' action stays disabled until both fields contain text, which prevents the
          most common validation error before the request is sent.
    \item The spinner replaces the icon of the action while the moderation proxy and the API are
          working, preventing a duplicate submission.
    \item The dialog closes only after a successful response; on failure it stays open with the
          draft intact.
\end{itemize}

\fnhead{Function Details}

\sublabel{Input Data}
\begin{specdatatable}{Data Name}{Description}
\texttt{categoryId} & Identifier of the category in which the post is created; a positive
integer sent by the dialog from the current category screen.\\ \hline
\texttt{title} & Post title, required, at most 255 characters in the database.\\ \hline
\texttt{content} & Post body, required; it is rich-text HTML produced by the editor.\\ \hline
\texttt{moderation} & Verdict attached by the moderation proxy: \texttt{status},
\texttt{score} between 0 and 1 and \texttt{categories}. It is an internal field, never typed by
the member.\\ \hline
Access token & Bearer token forwarded by the BFF proxy; the API resolves the author from it, so
the client cannot choose the author.\\ \hline
\end{specdatatable}

\sublabel{Validation Rules}
\begin{itemize}[leftmargin=1.6em,topsep=3pt,itemsep=2pt,parsep=0pt,label=$\bullet$]
    \item The title and the content are required and must contain at least one character after
          trimming.
    \item \texttt{categoryId} is optional for the API but must be a positive integer when it is
          present; the dialog always sends the category of the current screen.
    \item The category must exist and must not be soft deleted, otherwise the request fails with
          \texttt{Error.ForumCategoryNotFound}.
    \item The author is taken from the access token, never from the request body, so a member
          cannot create a post on behalf of somebody else.
    \item The moderation verdict must respect its own contract: a status among
          \texttt{PENDING}, \texttt{APPROVED} and \texttt{REJECTED}, and an optional score between
          0 and 1.
    \item The generated slug must remain unique; the timestamp and random suffix guarantee that two
          posts with the same title never collide.
\end{itemize}

\sublabel{Business Rules}
\begin{itemize}[leftmargin=1.6em,topsep=3pt,itemsep=2pt,parsep=0pt,label=$\bullet$]\brrulelist
    \item BR-FORUM-12: Only authenticated users can create posts.
    \item BR-FORUM-13: Every post must belong to exactly one user and may belong to a category.
    \item BR-FORUM-14: Title and content are required fields.
    \item BR-FORUM-15: A new post shall default to \texttt{moderationStatus = PENDING} unless
          moderation data is provided.
\end{itemize}

\sublabel{Normal Flow}
\begin{enumerate}[leftmargin=1.8em,topsep=3pt,itemsep=2pt,parsep=0pt]
    \item The member opens a category screen and selects ``New Post''.
    \item The client opens the dialog with the title of the target category.
    \item The member writes a title and a content; the ``Post'' action becomes available.
    \item The client trims both fields and sends the request through the BFF proxy.
    \item The proxy routes the request to the moderation proxy, which scans the text.
    \item The moderation proxy attaches its verdict and forwards the request to the REST API.
    \item The API validates the body and checks that the category exists.
    \item The API generates the slug, resolves the moderation status and inserts the post.
    \item The API returns the created post to the client.
    \item The client invalidates the cached lists of the category.
    \item When the status is \texttt{APPROVED}, the client closes the dialog and navigates to the
          new post.
    \item When the status is \texttt{PENDING}, the client closes the dialog, shows the message
          ``Your post has been saved and is awaiting moderator review.'' and returns the member to
          the list.
\end{enumerate}

\sublabel{Abnormal Flow / Exception Handling}
\begin{flowtable}{Case}{Description}
A1 & The title or the content is empty after trimming. The action is disabled, and a forced
request is rejected by the validation pipe with \texttt{Error.ForumPostTitleRequired} or
\texttt{Error.ForumPostContentRequired}.\\ \hline
A2 & The category was deleted while the dialog was open. The API answers
\texttt{Error.ForumCategoryNotFound} and the dialog stays open with the draft so the member can
choose another category.\\ \hline
A3 & The member is not signed in or the session expired. The BFF proxy refreshes the token once;
if the refresh fails the API rejects the request with an authentication error and the client
redirects to the sign-in page.\\ \hline
A4 & The moderation proxy is unavailable. The request fails before reaching the API, no post is
created, and the client shows the generic error message.\\ \hline
A5 & The moderation proxy flags the content. The post is stored with status \texttt{PENDING} or
\texttt{REJECTED}, is not publicly visible, and the member is informed that the post awaits
review.\\ \hline
A6 & The moderation proxy answers that the request must be reviewed asynchronously. The client
treats the response as a failure notification, shows the returned message and keeps the list
unchanged.\\ \hline
A7 & The insert fails at database level. The API returns \texttt{Error.FailedToCreateForumPost}
and the client keeps the dialog open so nothing written by the member is lost.\\ \hline
A8 & The member double clicks the action. The action is disabled while the mutation is pending,
therefore only one request is sent.\\ \hline
A9 & The member chooses ``Cancel''. The dialog closes, the draft is cleared and no request is
sent.\\ \hline
\end{flowtable}

\sublabel{Output Data}
\begin{specdatatable}{Data Name}{Description}
\texttt{id} & Identifier of the created post, used by the client to build the detail route.\\ \hline
\texttt{slug} & Generated slug, unique for every post.\\ \hline
\texttt{categoryId}, \texttt{userId} & Category and author of the post.\\ \hline
\texttt{title}, \texttt{content} & Stored title and stored rich-text content.\\ \hline
\texttt{moderationStatus} & \texttt{APPROVED}, \texttt{PENDING} or \texttt{REJECTED}; it decides
whether the post is immediately reachable.\\ \hline
\texttt{moderationScore}, \texttt{moderationCategories} & Verdict details kept for the moderation
queue; the score may be null when the proxy did not compute one.\\ \hline
\texttt{createdAt}, \texttt{updatedAt} & Timestamps of the insert.\\ \hline
\end{specdatatable}

\sublabel{Post-Conditions}
\begin{itemize}[leftmargin=1.6em,topsep=3pt,itemsep=2pt,parsep=0pt,label=$\bullet$]
    \item One post is inserted and belongs to the signed-in member and to the selected category.
    \item The post carries a moderation status and a unique slug.
    \item An approved post becomes visible in the category list immediately after the cache is
          invalidated; a pending or rejected post stays hidden until an administrator resolves it.
    \item The post appears in the moderation queue of the administrator.
\end{itemize}

% =================================================================

\subsubsection{Function: Edit Post}

\fnhead{Function Trigger}
The author of a post opens the post detail screen and selects the ``Edit'' action placed next to
the title. The action is rendered only for the author, therefore other members never see it.\newline
\textbf{Navigation path:} Home $\rightarrow$ Forum $\rightarrow$ \textit{Category} $\rightarrow$
\textit{Post} $\rightarrow$ Edit.

\fnhead{Function Description}

\sublabel{Actors / Roles}
\begin{itemize}[leftmargin=1.6em,topsep=3pt,itemsep=2pt,parsep=0pt,label=$\bullet$]
    \item \textbf{Primary Actor:} Freelancer or Client who authored the post.
    \item \textbf{Supporting Systems:} Frevia Web application, Frevia BFF proxy, Frevia REST API
          and PostgreSQL through Prisma ORM.
\end{itemize}

\sublabel{Purpose}
\noindent The function allows the author to correct the title or the body of a discussion without
losing the discussion itself, its comments and its likes. The post keeps its identity, its
creation date and its engagement, and the screen marks it as edited so readers know that the
content changed after publication.

\sublabel{Interface}
\begin{itemize}[leftmargin=1.6em,topsep=3pt,itemsep=2pt,parsep=0pt,label=$\bullet$]
    \item Modal dialog titled ``Edit Post'' with the description ``Update your post title and
          content.''
    \item Title input and rich-text editor pre-filled with the current values of the post.
    \item ``Cancel'' action, which leaves the post unchanged, and ``Save Changes'' action, which
          submits the new values.
    \item The post detail screen refreshes the title and the body in place after a successful
          save, and the author line starts showing ``Edited \{date\}''.
\end{itemize}

\sublabel{Data Processing Overview}
\begin{itemize}[leftmargin=1.6em,topsep=3pt,itemsep=2pt,parsep=0pt,label=$\bullet$]
    \item The client pre-fills the dialog from the post already loaded in the detail screen, so
          no additional read is required to open it.
    \item The client sends \texttt{PATCH /api/forums/posts/\{id\}} with the new title and content.
    \item The BFF proxy forwards the request directly to the REST API; only post creation goes
          through the moderation proxy, therefore an edition is not re-screened.
    \item The API loads the post with \texttt{deletedAt} null and rejects the request when the
          post does not exist.
    \item The API compares the author of the post with the identifier carried by the access token
          and rejects the request when they differ.
    \item The API updates the title, the content and the slug, then returns the updated post.
    \item The client updates the cached detail and the post lists of the category so the new title
          appears everywhere without a reload.
\end{itemize}

\fnhead{Screen Layout (Mock-up Description)}
\screenfigure[0.85]{image/screenshot/sprint1/kvy/EditPost.png}{Edit post dialog with the current values}
\begin{itemize}[leftmargin=1.6em,topsep=3pt,itemsep=2pt,parsep=0pt,label=$\bullet$]
    \item The dialog reuses the layout of the creation dialog, which keeps the editing experience
          consistent with publishing.
    \item Both fields are pre-filled, so the member edits the existing text instead of rewriting
          it.
    \item The ``Save Changes'' action stays disabled while a field is empty, mirroring the
          requirement that both fields remain mandatory.
    \item A spinner is shown during the request; the dialog closes only on success.
    \item After the save, the detail screen shows the new title, the new body and the edition
          marker next to the publication date.
\end{itemize}

\fnhead{Function Details}

\sublabel{Input Data}
\begin{specdatatable}{Data Name}{Description}
\texttt{id} & Identifier of the post taken from the current route.\\ \hline
\texttt{title} & New title, required, at most 255 characters in the database.\\ \hline
\texttt{content} & New body, required; rich-text HTML produced by the editor.\\ \hline
Access token & Identifies the requester; the API compares its identifier with the author of the
post.\\ \hline
\end{specdatatable}

\sublabel{Validation Rules}
\begin{itemize}[leftmargin=1.6em,topsep=3pt,itemsep=2pt,parsep=0pt,label=$\bullet$]
    \item The post must exist and must not be soft deleted; otherwise the request fails with
          \texttt{Error.ForumPostNotFound}.
    \item The requester must be the author of the post; otherwise the request fails with
          \texttt{Error.ForumPostNotOwned}. The check is performed on the server, so hiding the
          action in the interface is only a convenience.
    \item The title and the content are required; an empty value is rejected by the validation
          pipe before the update is executed.
    \item The slug is regenerated from the new title on every update, which keeps the URL
          readable and unique while the post identifier stays unchanged.
    \item The category of the post is preserved unless a new category identifier is supplied; the
          dialog never changes it.
\end{itemize}

\sublabel{Business Rules}
\begin{itemize}[leftmargin=1.6em,topsep=3pt,itemsep=2pt,parsep=0pt,label=$\bullet$]\brrulelist
    \item BR-FORUM-16: Only the author of a post can edit that post.
    \item BR-FORUM-17: Title and content are required fields.
    \item BR-FORUM-18: The original post identifier remains unchanged after editing; the slug is
          regenerated when the title changes.
\end{itemize}

\sublabel{Normal Flow}
\begin{enumerate}[leftmargin=1.8em,topsep=3pt,itemsep=2pt,parsep=0pt]
    \item The author opens the detail of a post they wrote.
    \item The client renders the ``Edit'' and ``Delete'' actions next to the title.
    \item The author selects ``Edit'' and the dialog opens with the current title and content.
    \item The author changes one or both fields.
    \item The client sends the update request with the access token.
    \item The API loads the post and confirms that it exists and is not deleted.
    \item The API confirms that the requester is the author.
    \item The API updates the title, the content and the slug.
    \item The API returns the updated post.
    \item The client closes the dialog, refreshes the detail and invalidates the category lists.
\end{enumerate}

\sublabel{Abnormal Flow / Exception Handling}
\begin{flowtable}{Case}{Description}
A1 & The post no longer exists or was deleted in the meantime. The API returns
\texttt{Error.ForumPostNotFound} and the client shows the error and closes the dialog.\\ \hline
A2 & The requester is not the author. The API returns \texttt{Error.ForumPostNotOwned}, the update
is refused and the client displays the returned message.\\ \hline
A3 & The title or the content is cleared before saving. The action is disabled, and a forced
request is rejected by the validation pipe.\\ \hline
A4 & The update fails at database level. The API returns
\texttt{Error.FailedToUpdateForumPost} and the dialog stays open with the edited values.\\ \hline
A5 & The author cancels the dialog. No request is sent and the post keeps its previous values.\\ \hline
A6 & The session expires while the dialog is open. The BFF proxy refreshes the token once; when
the refresh fails the request is rejected and the member is redirected to sign in.\\ \hline
A7 & Two editions are submitted in a row. The action is disabled while the mutation is pending,
so the last saved values are the ones returned by the server.\\ \hline
\end{flowtable}

\sublabel{Output Data}
\begin{specdatatable}{Data Name}{Description}
\texttt{id} & Identifier of the post, which never changes.\\ \hline
\texttt{title}, \texttt{content} & New values stored in the post.\\ \hline
\texttt{slug} & Slug regenerated from the new title.\\ \hline
\texttt{updatedAt} & New edition timestamp, used by the detail screen to display the ``Edited''
marker.\\ \hline
\texttt{createdAt} & Original publication timestamp, preserved.\\ \hline
\end{specdatatable}

\sublabel{Post-Conditions}
\begin{itemize}[leftmargin=1.6em,topsep=3pt,itemsep=2pt,parsep=0pt,label=$\bullet$]
    \item The post contains the new title and the new content and keeps its author, category,
          comments and likes.
    \item The detail screen and the category list display the new values.
    \item The post is marked as edited because \texttt{updatedAt} differs from \texttt{createdAt}.
    \item No other post and no comment is modified.
\end{itemize}

% =================================================================

\subsubsection{Function: Delete Post}

\fnhead{Function Trigger}
The author of a post, or an Administrator, selects the ``Delete'' action on the post detail screen
and confirms the operation in the confirmation dialog.\newline
\textbf{Navigation path:} Home $\rightarrow$ Forum $\rightarrow$ \textit{Category} $\rightarrow$
\textit{Post} $\rightarrow$ Delete.

\fnhead{Function Description}

\sublabel{Actors / Roles}
\begin{itemize}[leftmargin=1.6em,topsep=3pt,itemsep=2pt,parsep=0pt,label=$\bullet$]
    \item \textbf{Primary Actor:} the author of the post, or an Administrator acting on any post.
          The web interface renders the action for the author; the API also accepts an
          Administrator, who removes posts through the moderation and trash tools.
    \item \textbf{Supporting Systems:} Frevia Web application, Frevia BFF proxy, Frevia REST API
          and PostgreSQL through Prisma ORM.
\end{itemize}

\sublabel{Purpose}
\noindent The function removes a discussion from the public forum while preserving its record.
Deletion is soft: the row receives a deletion timestamp and is filtered out of every public
query, which keeps the data available for moderation, for the trash screen of the administrator
and for any later audit.

\sublabel{Interface}
\begin{itemize}[leftmargin=1.6em,topsep=3pt,itemsep=2pt,parsep=0pt,label=$\bullet$]
    \item ``Delete'' action with a bin icon, rendered for the author and for the Administrator.
    \item Confirmation dialog titled ``Delete Post'' with the warning ``Are you sure you want to
          delete this post? This action cannot be undone.''.
    \item ``Cancel'' action, which closes the dialog, and destructive ``Delete'' action, which
          performs the deletion.
    \item Redirect to the category screen after a successful deletion.
\end{itemize}

\sublabel{Data Processing Overview}
\begin{itemize}[leftmargin=1.6em,topsep=3pt,itemsep=2pt,parsep=0pt,label=$\bullet$]
    \item The client sends \texttt{DELETE /api/forums/posts/\{id\}} with the access token.
    \item The API loads the post with \texttt{deletedAt} null to confirm that it exists and is
          still visible.
    \item The API compares the author with the requester and also reads the role of the requester
          from the token.
    \item When the requester is neither the author nor an Administrator, the deletion is refused.
    \item When the check succeeds, the API writes the current date and time in the
          \texttt{deletedAt} column of the post.
    \item The client invalidates the post lists and navigates back to the category, where the post
          is no longer listed.
\end{itemize}

\fnhead{Screen Layout (Mock-up Description)}
\screenfigure[0.75]{image/screenshot/sprint1/kvy/DeletePost.png}{Delete post confirmation}
\begin{itemize}[leftmargin=1.6em,topsep=3pt,itemsep=2pt,parsep=0pt,label=$\bullet$]
    \item The confirmation dialog repeats the name of the object being deleted, which prevents
          deleting the wrong discussion.
    \item The warning states that the action cannot be undone from the member point of view,
          because the post disappears from the forum immediately.
    \item The destructive style of the confirmation action distinguishes it from the cancel
          action.
    \item After confirmation the screen returns to the category list, so the member sees the post
          removed from the discussion list.
\end{itemize}

\fnhead{Function Details}

\sublabel{Input Data}
\begin{specdatatable}{Data Name}{Description}
\texttt{id} & Identifier of the post to delete, taken from the confirmation dialog.\\ \hline
Access token & Identifies the requester and carries the role used to authorize an
Administrator.\\ \hline
\texttt{deletedAt} & Deletion timestamp written by the API; it is internal, not supplied by the
client.\\ \hline
\end{specdatatable}

\sublabel{Validation Rules}
\begin{itemize}[leftmargin=1.6em,topsep=3pt,itemsep=2pt,parsep=0pt,label=$\bullet$]
    \item The post must exist and must not be already deleted; otherwise the API returns
          \texttt{Error.ForumPostNotFound}.
    \item The requester must be the author of the post or must hold the Administrator role;
          otherwise the API returns \texttt{Error.ForumPostNotOwned}.
    \item The deletion timestamp is produced by the server, so the client cannot backdate or
          falsify it.
    \item Only one row is affected per request; the comments and the likes of the post are
          preserved in the database and simply stop being reachable through the post.
\end{itemize}

\sublabel{Business Rules}
\begin{itemize}[leftmargin=1.6em,topsep=3pt,itemsep=2pt,parsep=0pt,label=$\bullet$]\brrulelist
    \item BR-FORUM-19: Only the author of a post or an Admin can delete that post.
    \item BR-FORUM-20: Post deletion shall be implemented as a soft delete using the
          \texttt{deletedAt} field.
    \item BR-FORUM-21: Soft-deleted posts shall not appear in public forum listings.
\end{itemize}

\sublabel{Normal Flow}
\begin{enumerate}[leftmargin=1.8em,topsep=3pt,itemsep=2pt,parsep=0pt]
    \item The author or the Administrator opens the post detail screen.
    \item The client renders the ``Delete'' action, which is available for this role.
    \item The member selects the action and the confirmation dialog is displayed.
    \item The member confirms the deletion.
    \item The client sends the delete request with the access token.
    \item The API loads the post and confirms that it exists and is not deleted.
    \item The API confirms that the requester is the author or an Administrator.
    \item The API writes the deletion timestamp in the post row.
    \item The API returns the deleted post.
    \item The client invalidates the lists, closes the dialog and routes back to the category.
    \item The category list no longer contains the post and its counters are recomputed without
          it.
\end{enumerate}

\sublabel{Abnormal Flow / Exception Handling}
\begin{flowtable}{Case}{Description}
A1 & The post does not exist or was already deleted. The API returns
\texttt{Error.ForumPostNotFound} and the client shows the error message.\\ \hline
A2 & A member tries to delete a post written by somebody else. The API returns
\texttt{Error.ForumPostNotOwned} and nothing is deleted.\\ \hline
A3 & An Administrator deletes the post of another member. The operation is allowed and the post
leaves the public forum immediately.\\ \hline
A4 & The database update fails. The API returns \texttt{Error.FailedToDeleteForumPost} and the
post remains visible, so the member can retry.\\ \hline
A5 & The member cancels the confirmation. No request is sent and the post is unchanged.\\ \hline
A6 & The session expires before the confirmation. The request is rejected after the single token
refresh attempt and the member is redirected to sign in.\\ \hline
A7 & Another member opens the same post while it is being deleted. The next request returns
\texttt{Error.ForumPostNotFound} and their screen shows the not-found state.\\ \hline
\end{flowtable}

\sublabel{Output Data}
\begin{specdatatable}{Data Name}{Description}
\texttt{id} & Identifier of the deleted post.\\ \hline
\texttt{userId}, \texttt{categoryId} & Author and category of the deleted post, returned for
confirmation.\\ \hline
\texttt{title}, \texttt{slug} & Title and slug of the deleted post.\\ \hline
\texttt{deletedAt} & Server timestamp recorded in the row; the post is soft deleted from this
moment.\\ \hline
\end{specdatatable}

\sublabel{Post-Conditions}
\begin{itemize}[leftmargin=1.6em,topsep=3pt,itemsep=2pt,parsep=0pt,label=$\bullet$]
    \item The post is excluded from the category list, from the post detail and from the landing
          page counters.
    \item The row and its related comments and likes are retained for moderation and for the trash
          screen of the Administrator.
    \item The member is returned to the category screen, which reflects the new totals.
\end{itemize}

% =================================================================

\subsubsection{Function: Like a Post}

\fnhead{Function Trigger}
A signed-in member selects the heart button in the action bar of a post. A guest who selects the
button is routed to the sign-in page instead of producing a request.\newline
\textbf{Navigation path:} Home $\rightarrow$ Forum $\rightarrow$ \textit{Category} $\rightarrow$
\textit{Post} $\rightarrow$ Like.

\fnhead{Function Description}

\sublabel{Actors / Roles}
\begin{itemize}[leftmargin=1.6em,topsep=3pt,itemsep=2pt,parsep=0pt,label=$\bullet$]
    \item \textbf{Primary Actor:} Freelancer and Client. An Administrator can also react while
          browsing.
    \item \textbf{Supporting Systems:} Frevia Web application, Frevia BFF proxy, Frevia REST API
          and PostgreSQL through Prisma ORM.
\end{itemize}

\sublabel{Purpose}
\noindent The function lets a member express appreciation for a discussion and lets the community
measure which posts are relevant. It is implemented as a toggle: the same action adds the
reaction when the member has not liked the post yet and removes it when the member already liked
it, so a member can never like the same post twice.

\sublabel{Interface}
\begin{itemize}[leftmargin=1.6em,topsep=3pt,itemsep=2pt,parsep=0pt,label=$\bullet$]
    \item Heart button with the like counter, placed in the action bar under the post body.
    \item Outline heart for a post that the member has not liked, filled heart for a post that the
          member liked.
    \item The button is disabled while the request is running, which prevents repeated toggles
          from a single intention.
\end{itemize}

\sublabel{Data Processing Overview}
\begin{itemize}[leftmargin=1.6em,topsep=3pt,itemsep=2pt,parsep=0pt,label=$\bullet$]
    \item The client sends \texttt{POST /api/forums/posts/\{id\}/like} to the BFF proxy, which
          forwards it to the REST API.
    \item The API loads the post with \texttt{deletedAt} null and answers with
          \texttt{Error.ForumPostNotFound} when it is missing.
    \item The API looks for a like row that matches the member and the post.
    \item When the row exists, the API removes the reaction and answers with
          \texttt{liked: false}.
    \item When the row does not exist, the API inserts a like row with the reaction type
          \texttt{like} and answers with \texttt{liked: true}.
    \item The client invalidates the likes query of the post, which refreshes the counter and the
          filled state of the heart.
    \item The unique constraint of the like table on the pair member and post guarantees that a
          duplicate reaction cannot be stored, even if two requests are issued at the same time.
\end{itemize}

\fnhead{Screen Layout (Mock-up Description)}
\screenfigurepair{image/screenshot/sprint1/kvy/BeforeLikePost.png}{image/screenshot/sprint1/kvy/AfterLikePost.png}{Post before and after being liked}
\begin{itemize}[leftmargin=1.6em,topsep=3pt,itemsep=2pt,parsep=0pt,label=$\bullet$]
    \item Before the reaction the heart is outlined and the counter shows the current number of
          likes.
    \item After the reaction the heart is filled, coloured with the accent of the application and
          slightly enlarged, and the counter is incremented by one.
    \item The counter is the single source of truth for the likes list: the same number is shown
          on the detail screen, in the likes dialog and on the post cards of the category.
\end{itemize}

\fnhead{Function Details}

\sublabel{Input Data}
\begin{specdatatable}{Data Name}{Description}
\texttt{id} & Identifier of the post taken from the current route.\\ \hline
Access token & Identifies the member who reacts; the API derives the like owner from it.\\ \hline
Reaction type & Constant value \texttt{like} stored in the reactivity row; the interface does not
offer another reaction for the moment.\\ \hline
\end{specdatatable}

\sublabel{Validation Rules}
\begin{itemize}[leftmargin=1.6em,topsep=3pt,itemsep=2pt,parsep=0pt,label=$\bullet$]
    \item The requester must be authenticated; a request without a valid token is rejected with an
          authentication error.
    \item The post must exist and must not be soft deleted.
    \item A member can hold only one reaction per post, enforced by the unique constraint on the
          pair member and post.
    \item The reaction is always attached to the member of the token; the client cannot send a
          different member identifier.
\end{itemize}

\sublabel{Business Rules}
\begin{itemize}[leftmargin=1.6em,topsep=3pt,itemsep=2pt,parsep=0pt,label=$\bullet$]\brrulelist
    \item BR-FORUM-22: Only authenticated users can like posts.
    \item BR-FORUM-23: A user can have only one active like reaction per post.
    \item BR-FORUM-24: The reaction shall be stored in the \texttt{ForumLike} table.
    \item BR-FORUM-25 to BR-FORUM-27: removing a reaction is allowed only when the corresponding
          like record exists, and it deletes that record.
\end{itemize}

\sublabel{Normal Flow}
\begin{enumerate}[leftmargin=1.8em,topsep=3pt,itemsep=2pt,parsep=0pt]
    \item The signed-in member opens a post they have not liked yet.
    \item The client displays an outlined heart next to the current counter.
    \item The member selects the heart and the client sends the toggle request.
    \item The API confirms that the post is visible.
    \item The API finds no reaction of this member on this post.
    \item The API inserts the like row.
    \item The API answers with \texttt{liked: true}.
    \item The client refetches the likes, fills the heart and increments the counter.
    \item When the member selects the heart again, the same sequence removes the row and the
          interface returns to the outlined heart with the previous counter.
\end{enumerate}

\sublabel{Abnormal Flow / Exception Handling}
\begin{flowtable}{Case}{Description}
A1 & The post does not exist or was deleted. The API returns
\texttt{Error.ForumPostNotFound} and the counter is left unchanged.\\ \hline
A2 & The member is not signed in. The client routes to the sign-in page before any request is
sent, and a direct request is rejected with an authentication error.\\ \hline
A3 & The member already liked the post. The request removes the reaction instead of adding a
second one, which is the intended toggle behaviour.\\ \hline
A4 & The insert fails at database level. The API returns
\texttt{Error.FailedToLikeForumPost} and the interface keeps the previous state.\\ \hline
A5 & The member double clicks the heart. The button is disabled while the request runs, so only
one reaction is created.\\ \hline
A6 & The likes query cannot be refreshed after a successful toggle. The heart keeps the state
returned by the action and the counter is refreshed on the next visit.\\ \hline
\end{flowtable}

\sublabel{Output Data}
\begin{specdatatable}{Data Name}{Description}
\texttt{liked} & \texttt{true} when the reaction was added, \texttt{false} when it was
removed.\\ \hline
Like row & Member, post, reaction type and creation date of the stored reaction.\\ \hline
Updated counter & The like counter of the post, recomputed by the client from the refreshed likes
list.\\ \hline
\end{specdatatable}

\sublabel{Post-Conditions}
\begin{itemize}[leftmargin=1.6em,topsep=3pt,itemsep=2pt,parsep=0pt,label=$\bullet$]
    \item The post holds exactly one reaction of the member, or none when the action removed it.
    \item The counter displayed on the post matches the number of likes stored for that post.
    \item The member appears in, or disappears from, the likes list of the post.
\end{itemize}

% =================================================================

\subsubsection{Function: View Post Likes}

\fnhead{Function Trigger}
A visitor selects the likes counter under the post body, which opens the list of the members who
liked the post. The function is public, therefore it works for guests as well.\newline
\textbf{Navigation path:} Home $\rightarrow$ Forum $\rightarrow$ \textit{Category} $\rightarrow$
\textit{Post} $\rightarrow$ Likes.

\fnhead{Function Description}

\sublabel{Actors / Roles}
\begin{itemize}[leftmargin=1.6em,topsep=3pt,itemsep=2pt,parsep=0pt,label=$\bullet$]
    \item \textbf{Primary Actor:} Guest, Freelancer, Client and Administrator.
    \item \textbf{Supporting Systems:} Frevia Web application, Frevia BFF proxy, Frevia REST API
          and PostgreSQL through Prisma ORM.
\end{itemize}

\sublabel{Purpose}
\noindent The function answers the question ``who liked this post''. It also feeds the state of
the like button, because the detail screen derives both the counter and the liked state of the
current member from the same list.

\sublabel{Interface}
\begin{itemize}[leftmargin=1.6em,topsep=3pt,itemsep=2pt,parsep=0pt,label=$\bullet$]
    \item Dialog titled ``Likes'' with the description ``People who liked this post.''.
    \item Row per member with the avatar, the display name and, when the profile is incomplete, a
          fallback name built from the member identifier.
    \item Empty message ``No likes yet.'' when the post has not been liked.
\end{itemize}

\sublabel{Data Processing Overview}
\begin{itemize}[leftmargin=1.6em,topsep=3pt,itemsep=2pt,parsep=0pt,label=$\bullet$]
    \item The client calls \texttt{GET /api/forums/posts/\{id\}/like} and caches the result under a
          key that contains the post identifier.
    \item The API reads every like row of the post and joins the profile of the member.
    \item The payload contains the identifier, the member, the reaction type and the creation date
          of the reaction.
    \item The detail screen uses the same payload to count the likes and to detect the reaction of
          the signed-in member by comparing identifiers.
    \item The dialog renders the rows in the order returned by the query, which follows the
          insertion order of the reactions.
\end{itemize}

\fnhead{Screen Layout (Mock-up Description)}
\screenfigure[0.6]{image/screenshot/sprint1/kvy/ViewPostLike.png}{Likes dialog listing the members who liked the post}
\begin{itemize}[leftmargin=1.6em,topsep=3pt,itemsep=2pt,parsep=0pt,label=$\bullet$]
    \item The likes counter placed under the post body opens the dialog and states how many
          members liked the discussion.
    \item Each row shows the avatar and the display name of a member, which links the reaction to
          a visible identity of the community.
    \item When nobody liked the post, the dialog states ``No likes yet.'' instead of showing an
          empty panel.
\end{itemize}

\fnhead{Function Details}

\sublabel{Input Data}
\begin{specdatatable}{Data Name}{Description}
\texttt{id} & Identifier of the post whose reactions are listed.\\ \hline
Session identity & Optional; used by the detail screen to highlight the role of the current member
in the list, and not needed by the endpoint itself.\\ \hline
\end{specdatatable}

\sublabel{Validation Rules}
\begin{itemize}[leftmargin=1.6em,topsep=3pt,itemsep=2pt,parsep=0pt,label=$\bullet$]
    \item The post identifier must be a positive integer accepted by the route parameter pipe.
    \item The endpoint is public, so no token is required and no authorization rule applies.
    \item A post without reaction returns an empty array, which is a valid answer and not an
          error.
    \item Profile fields may be null; the interface then substitutes the fallback name and the
          initial of the display name for the avatar.
\end{itemize}

\sublabel{Business Rules}
\begin{itemize}[leftmargin=1.6em,topsep=3pt,itemsep=2pt,parsep=0pt,label=$\bullet$]\brrulelist
    \item BR-FORUM-28: The like detail shall include the user identifier and the user profile
          information for each liker.
    \item BR-FORUM-23: A user can have only one active like reaction per post, therefore a member
          never appears twice in the list.
\end{itemize}

\sublabel{Normal Flow}
\begin{enumerate}[leftmargin=1.8em,topsep=3pt,itemsep=2pt,parsep=0pt]
    \item The visitor opens a post detail screen.
    \item The client requests the likes of the post to compute the counter.
    \item The API returns every reaction of the post with the profile of the member.
    \item The client shows the counter under the post body.
    \item The visitor selects the counter and the dialog opens with the list of members.
    \item The visitor closes the dialog and returns to the discussion.
\end{enumerate}

\sublabel{Abnormal Flow / Exception Handling}
\begin{flowtable}{Case}{Description}
A1 & The post has no reaction. The dialog shows ``No likes yet.'' and the counter stays at
zero.\\ \hline
A2 & A member has no profile record. The row shows the fallback name ``User \#\{id\}'' and the
initial letter in place of the avatar.\\ \hline
A3 & The request fails. The API returns \texttt{Error.FailedToViewForumLikeDetail}, the dialog
shows an error message and the counter keeps its last known value.\\ \hline
A4 & The post was deleted while the dialog was open. The list refresh returns no row and the
dialog shows the empty state.\\ \hline
A5 & The visitor is offline. The dialog shows the generic error message and can be closed, which
leaves the rest of the screen usable.\\ \hline
\end{flowtable}

\sublabel{Output Data}
\begin{specdatatable}{Data Name}{Description}
\texttt{likes[]} & One entry per reaction with \texttt{id}, \texttt{userId}, \texttt{postId},
\texttt{commentId}, \texttt{reactionType} and \texttt{createdAt}.\\ \hline
\texttt{user} & Identifier, display name and avatar of the member who reacted.\\ \hline
Derived count & Length of the list, displayed as the likes counter of the post.\\ \hline
\end{specdatatable}

\sublabel{Post-Conditions}
\begin{itemize}[leftmargin=1.6em,topsep=3pt,itemsep=2pt,parsep=0pt,label=$\bullet$]
    \item The visitor sees every member who liked the post, without duplicates.
    \item The like counter of the post matches the number of rows returned.
    \item No data is modified: the function is read-only.
\end{itemize}

% =================================================================

\subsubsection{Function: Create Comment}

\fnhead{Function Trigger}
A signed-in member writes in the comment box placed under the post body and submits it with the
``Comment'' action or with the \texttt{Ctrl+Enter} shortcut. A guest sees an invitation to sign in
instead of the editor.\newline
\textbf{Navigation path:} Home $\rightarrow$ Forum $\rightarrow$ \textit{Category} $\rightarrow$
\textit{Post} $\rightarrow$ Comments.

\fnhead{Function Description}

\sublabel{Actors / Roles}
\begin{itemize}[leftmargin=1.6em,topsep=3pt,itemsep=2pt,parsep=0pt,label=$\bullet$]
    \item \textbf{Primary Actor:} Freelancer and Client. An Administrator can also comment while
          browsing the forum.
    \item \textbf{Supporting Systems:} Frevia Web application, Frevia BFF proxy, Frevia REST API
          and PostgreSQL through Prisma ORM.
\end{itemize}

\sublabel{Purpose}
\noindent The function lets a member answer a discussion, ask a question or share an experience
inside the thread of a post. Comments are attached to one post and one author, can be edited and
deleted by their owner, and carry their own like counter.

\sublabel{Interface}
\begin{itemize}[leftmargin=1.6em,topsep=3pt,itemsep=2pt,parsep=0pt,label=$\bullet$]
    \item Comment card between the post body and the thread, with the avatar of the member and a
          text area using the placeholder ``Write a comment\ldots''.
    \item Hint ``Press \texttt{Ctrl+Enter} to submit'' below the text area.
    \item ``Comment'' action with a send icon, disabled while the field is empty or a request is
          running.
    \item Guest state: the editor is replaced by the sentence ``Log in to join the discussion.''
          with a link to the sign-in page.
\end{itemize}

\sublabel{Data Processing Overview}
\begin{itemize}[leftmargin=1.6em,topsep=3pt,itemsep=2pt,parsep=0pt,label=$\bullet$]
    \item The client trims the text and sends \texttt{POST /api/forums/posts/\{postId\}/comments}
          with the content.
    \item The BFF proxy forwards the request to the REST API, which validates the body.
    \item The API verifies that the post exists and is not soft deleted.
    \item The API inserts the comment with the author taken from the access token.
    \item The response contains the new comment with its author profile, a like counter of zero and
          an unliked state, because a comment cannot be liked before it exists.
    \item The client clears the editor and invalidates the comment queries of the post, which
          refreshes the thread and the comment counter.
\end{itemize}

\fnhead{Screen Layout (Mock-up Description)}
\screenfigure[0.85]{image/screenshot/sprint1/kvy/CreateComment.png}{Comment editor under the post body}
\begin{itemize}[leftmargin=1.6em,topsep=3pt,itemsep=2pt,parsep=0pt,label=$\bullet$]
    \item The editor sits directly under the action bar, so the member writes without leaving the
          discussion.
    \item The counter next to the ``Comments'' heading shows how many comments the post currently
          holds.
    \item The ``Comment'' action stays disabled while the field is empty, which prevents empty
          submissions.
    \item The hint about \texttt{Ctrl+Enter} makes the keyboard shortcut discoverable.
    \item Every submitted comment appears at the top of the thread, with the avatar, the display
          name, the date and the like and edit actions.
\end{itemize}

\fnhead{Function Details}

\sublabel{Input Data}
\begin{specdatatable}{Data Name}{Description}
\texttt{postId} & Identifier of the post taken from the route; the comment is attached to it.\\ \hline
\texttt{content} & Comment text, required, at least one character after trimming.\\ \hline
Access token & Identifies the author of the comment; it is read by the API and never sent in the
body.\\ \hline
\end{specdatatable}

\sublabel{Validation Rules}
\begin{itemize}[leftmargin=1.6em,topsep=3pt,itemsep=2pt,parsep=0pt,label=$\bullet$]
    \item The requester must be authenticated; the API rejects a request without a valid token.
    \item The content must not be empty; the validation pipe rejects an empty or whitespace-only
          body.
    \item The post must exist and must not be soft deleted, otherwise the request fails with
          \texttt{Error.ForumPostNotFound}.
    \item The author of the comment is always the member of the token, so a member cannot post a
          comment on behalf of somebody else.
\end{itemize}

\sublabel{Business Rules}
\begin{itemize}[leftmargin=1.6em,topsep=3pt,itemsep=2pt,parsep=0pt,label=$\bullet$]\brrulelist
    \item BR-FORUM-32: Only authenticated users can create comments.
    \item BR-FORUM-33: Comment content cannot be empty.
    \item BR-FORUM-34: Each comment shall be associated with a valid post and user.
\end{itemize}

\sublabel{Normal Flow}
\begin{enumerate}[leftmargin=1.8em,topsep=3pt,itemsep=2pt,parsep=0pt]
    \item The signed-in member opens a post and scrolls to the comment section.
    \item The client renders the editor with the placeholder and the hint.
    \item The member writes the comment and submits it with the action or the shortcut.
    \item The client trims the text and sends the request.
    \item The API validates the content and loads the post.
    \item The API inserts the comment with the author of the session.
    \item The API returns the created comment with an empty like state.
    \item The client clears the editor and refreshes the thread.
    \item The new comment is displayed at the top of the list and the counter is incremented.
\end{enumerate}

\sublabel{Abnormal Flow / Exception Handling}
\begin{flowtable}{Case}{Description}
A1 & The content is empty. The action is disabled, and a forced request is rejected by the
validation pipe without touching the database.\\ \hline
A2 & The post does not exist or was deleted. The API returns \texttt{Error.ForumPostNotFound} and
the client displays the error message while keeping the draft.\\ \hline
A3 & The member is not signed in. The editor is replaced by the sign-in invitation, and a direct
request is rejected with an authentication error.\\ \hline
A4 & The insert fails at database level. The API returns
\texttt{Error.FailedToCreateForumComment} and the editor keeps the text so it is not lost.\\ \hline
A5 & The member submits twice in a row. The action is disabled while the request runs, which
prevents duplicated comments.\\ \hline
A6 & The session expires during the submission. The BFF proxy attempts one refresh; when it fails
the request is rejected and the member is redirected to sign in.\\ \hline
\end{flowtable}

\sublabel{Output Data}
\begin{specdatatable}{Data Name}{Description}
\texttt{id} & Identifier of the created comment.\\ \hline
\texttt{postId} & Post the comment belongs to.\\ \hline
\texttt{content} & Stored comment text.\\ \hline
\texttt{user} & Identifier, display name and avatar of the author.\\ \hline
\texttt{likeCount}, \texttt{likedByMe} & Zero and false for a new comment.\\ \hline
\texttt{createdAt}, \texttt{updatedAt} & Timestamps of the insert.\\ \hline
\end{specdatatable}

\sublabel{Post-Conditions}
\begin{itemize}[leftmargin=1.6em,topsep=3pt,itemsep=2pt,parsep=0pt,label=$\bullet$]
    \item One comment is stored and belongs to the signed-in member and to the post.
    \item The thread and the comment counter of the post reflect the new comment.
    \item The comment is immediately visible to every visitor of the post.
\end{itemize}

% =================================================================

\subsubsection{Function: Update Comment}

\fnhead{Function Trigger}
The author of a comment selects the pencil action next to their comment, edits the text and
confirms the change. The action is rendered only for the author of the comment.\newline
\textbf{Navigation path:} Home $\rightarrow$ Forum $\rightarrow$ \textit{Category} $\rightarrow$
\textit{Post} $\rightarrow$ Comments $\rightarrow$ Edit.

\fnhead{Function Description}

\sublabel{Actors / Roles}
\begin{itemize}[leftmargin=1.6em,topsep=3pt,itemsep=2pt,parsep=0pt,label=$\bullet$]
    \item \textbf{Primary Actor:} Freelancer or Client who authored the comment.
    \item \textbf{Supporting Systems:} Frevia Web application, Frevia BFF proxy, Frevia REST API
          and PostgreSQL through Prisma ORM.
\end{itemize}

\sublabel{Purpose}
\noindent The function lets a member correct a comment after publication, for example to fix a
typo or to add a clarification, without deleting the reply and losing its position in the
thread. The comment keeps its identifier, its author and its likes, and the new text replaces the
previous one.

\sublabel{Interface}
\begin{itemize}[leftmargin=1.6em,topsep=3pt,itemsep=2pt,parsep=0pt,label=$\bullet$]
    \item Pencil action on the header row of the comment, next to the author and the date.
    \item Inline editor pre-filled with the current text of the comment.
    \item ``Save'' action to confirm and ``Cancel'' action to discard the change.
    \item The thread is refreshed after the save, and the updated comment is displayed with its
          new text.
\end{itemize}

\sublabel{Data Processing Overview}
\begin{itemize}[leftmargin=1.6em,topsep=3pt,itemsep=2pt,parsep=0pt,label=$\bullet$]
    \item The client sends \texttt{PATCH /api/forums/posts/\{postId\}/comments/\{commentId\}} with
          the new content.
    \item The API loads the comment with \texttt{deletedAt} null and rejects the request when it
          does not exist.
    \item The API verifies that the comment belongs to the post of the route, so a comment cannot
          be edited through the route of another post.
    \item The API compares the author of the comment with the member of the token and refuses the
          request when they differ.
    \item The API updates the content and returns the comment with its author profile.
    \item The client invalidates the comment queries of the post and re-renders the thread.
\end{itemize}

\fnhead{Screen Layout (Mock-up Description)}
\screenfigure[0.85]{image/screenshot/sprint1/kvy/UpdateComment.png}{Inline editor of an existing comment}
\begin{itemize}[leftmargin=1.6em,topsep=3pt,itemsep=2pt,parsep=0pt,label=$\bullet$]
    \item The editor replaces the comment body in place, so the member keeps the context of the
          thread while editing.
    \item The existing text is pre-filled, which makes small corrections fast.
    \item The save action is disabled while the text is empty, because the content stays a
          required field.
    \item After the save the comment shows the edited text and the thread keeps the same order.
\end{itemize}

\fnhead{Function Details}

\sublabel{Input Data}
\begin{specdatatable}{Data Name}{Description}
\texttt{postId} & Identifier of the post from the route; it scopes the comment.\\ \hline
\texttt{commentId} & Identifier of the comment to update.\\ \hline
\texttt{content} & New comment text, required, at least one character.\\ \hline
Access token & Identifies the requester; the API compares it with the comment author.\\ \hline
\end{specdatatable}

\sublabel{Validation Rules}
\begin{itemize}[leftmargin=1.6em,topsep=3pt,itemsep=2pt,parsep=0pt,label=$\bullet$]
    \item The comment must exist and must not be soft deleted; otherwise the request fails with
          \texttt{Error.ForumCommentNotFound}.
    \item The comment must belong to the post of the route; a mismatch is reported as a missing
          comment, which hides the existence of the comment through another route.
    \item The requester must be the author of the comment; otherwise the request fails with
          \texttt{Error.ForumCommentNotOwner}.
    \item The content is required and must not be empty after trimming.
    \item The update touches the content only: the author, the post, the creation date and the
          likes of the comment are preserved.
\end{itemize}

\sublabel{Business Rules}
\begin{itemize}[leftmargin=1.6em,topsep=3pt,itemsep=2pt,parsep=0pt,label=$\bullet$]\brrulelist
    \item BR-FORUM-35: Only the owner of a comment can edit it.
    \item BR-FORUM-36: Comment content cannot be empty.
    \item BR-FORUM-37: A comment must belong to the specified post to be edited.
\end{itemize}

\sublabel{Normal Flow}
\begin{enumerate}[leftmargin=1.8em,topsep=3pt,itemsep=2pt,parsep=0pt]
    \item The author of a comment selects the pencil action on their own comment.
    \item The client opens the inline editor with the current text.
    \item The author modifies the text and confirms.
    \item The client sends the update request with the access token.
    \item The API loads the comment and confirms that it exists.
    \item The API confirms that the comment belongs to the post of the route.
    \item The API confirms that the requester is the author.
    \item The API stores the new content.
    \item The API returns the updated comment.
    \item The client refreshes the thread and displays the new text.
\end{enumerate}

\sublabel{Abnormal Flow / Exception Handling}
\begin{flowtable}{Case}{Description}
A1 & The comment does not exist or was deleted. The API returns
\texttt{Error.ForumCommentNotFound} and the editor closes with an error message.\\ \hline
A2 & The comment belongs to another post. The request is refused with the same not-found error,
which prevents editing a comment through a foreign route.\\ \hline
A3 & The requester is not the author. The API returns \texttt{Error.ForumCommentNotOwner} and the
text is not changed.\\ \hline
A4 & The new text is empty. The save action is disabled, and a forced request is rejected by the
validation pipe.\\ \hline
A5 & The update fails at database level. The API returns
\texttt{Error.FailedToUpdateForumComment} and the editor keeps the text.\\ \hline
A6 & The author cancels the edition. No request is sent and the previous text is restored.\\ \hline
\end{flowtable}

\sublabel{Output Data}
\begin{specdatatable}{Data Name}{Description}
\texttt{id} & Identifier of the updated comment.\\ \hline
\texttt{postId} & Post the comment belongs to, unchanged.\\ \hline
\texttt{content} & New comment text.\\ \hline
\texttt{user} & Identifier and profile of the author, unchanged.\\ \hline
\texttt{updatedAt} & New edition timestamp.\\ \hline
\texttt{createdAt} & Original creation timestamp, preserved.\\ \hline
\end{specdatatable}

\sublabel{Post-Conditions}
\begin{itemize}[leftmargin=1.6em,topsep=3pt,itemsep=2pt,parsep=0pt,label=$\bullet$]
    \item The comment holds the new text and keeps its position, its author and its likes.
    \item The thread displays the updated comment after the refresh.
    \item No other comment and no post is modified.
\end{itemize}

% =================================================================

\subsubsection{Function: Delete Comment}

\fnhead{Function Trigger}
The author of a comment, or an Administrator, selects the bin action next to a comment and
confirms the deletion in the confirmation dialog.\newline
\textbf{Navigation path:} Home $\rightarrow$ Forum $\rightarrow$ \textit{Category} $\rightarrow$
\textit{Post} $\rightarrow$ Comments $\rightarrow$ Delete.

\fnhead{Function Description}

\sublabel{Actors / Roles}
\begin{itemize}[leftmargin=1.6em,topsep=3pt,itemsep=2pt,parsep=0pt,label=$\bullet$]
    \item \textbf{Primary Actor:} the author of the comment, or an Administrator acting on any
          comment. The web interface exposes the action to the author of the comment; the API also
          accepts an Administrator, who deletes comments through the moderation tools.
    \item \textbf{Supporting Systems:} Frevia Web application, Frevia BFF proxy, Frevia REST API
          and PostgreSQL through Prisma ORM.
\end{itemize}

\sublabel{Purpose}
\noindent The function removes a reply from the thread. As for posts, the deletion is soft: the
comment receives a deletion timestamp and is filtered out of every public query, so the record
remains available for moderation and for the trash screen of the Administrator while the
discussion stays coherent for the readers.

\sublabel{Interface}
\begin{itemize}[leftmargin=1.6em,topsep=3pt,itemsep=2pt,parsep=0pt,label=$\bullet$]
    \item Bin action on the header row of the comment, next to the edit action.
    \item Confirmation dialog titled ``Delete Comment'' with the warning ``Are you sure you want
          to delete this comment? This action cannot be undone.''.
    \item ``Cancel'' action and destructive ``Delete'' action.
    \item The thread is refreshed after the deletion and the comment counter decreases.
\end{itemize}

\sublabel{Data Processing Overview}
\begin{itemize}[leftmargin=1.6em,topsep=3pt,itemsep=2pt,parsep=0pt,label=$\bullet$]
    \item The client sends \texttt{DELETE /api/forums/posts/\{postId\}/comments/\{commentId\}}.
    \item The API loads the comment with \texttt{deletedAt} null and rejects the request when it is
          missing.
    \item The API verifies that the comment belongs to the post of the route.
    \item The API accepts the requester when it is the author or when the token carries the
          Administrator role.
    \item The API writes the deletion timestamp in the comment row.
    \item The client invalidates the comment queries and refreshes the thread and the counter.
\end{itemize}

\fnhead{Screen Layout (Mock-up Description)}
\screenfigure[0.8]{image/screenshot/sprint1/kvy/DeleteComment.png}{Delete comment confirmation}
\begin{itemize}[leftmargin=1.6em,topsep=3pt,itemsep=2pt,parsep=0pt,label=$\bullet$]
    \item The confirmation dialog names the object, so a member who owns several comments cannot
          delete the wrong one.
    \item The destructive style of the confirm action separates it from the cancel action and
          warns about the consequence.
    \item After the confirmation the comment disappears from the thread and the comment counter of
          the post is updated.
    \item The remaining comments keep their order, because the thread is sorted by creation date.
\end{itemize}

\fnhead{Function Details}

\sublabel{Input Data}
\begin{specdatatable}{Data Name}{Description}
\texttt{postId} & Identifier of the post from the route; it scopes the comment.\\ \hline
\texttt{commentId} & Identifier of the comment to delete.\\ \hline
Access token & Identifies the requester and carries the role used to authorize an
Administrator.\\ \hline
\texttt{deletedAt} & Deletion timestamp written by the server, not supplied by the client.\\ \hline
\end{specdatatable}

\sublabel{Validation Rules}
\begin{itemize}[leftmargin=1.6em,topsep=3pt,itemsep=2pt,parsep=0pt,label=$\bullet$]
    \item The comment must exist and must not be already deleted; otherwise the request fails with
          \texttt{Error.ForumCommentNotFound}.
    \item The comment must belong to the post of the route; a mismatch is reported as a missing
          comment.
    \item The requester must be the author of the comment or must hold the Administrator role;
          otherwise the request fails with \texttt{Error.ForumCommentNotOwner}.
    \item The deletion timestamp is produced by the server.
    \item The post itself is not deleted, and the likes of the comment are kept in the database.
\end{itemize}

\sublabel{Business Rules}
\begin{itemize}[leftmargin=1.6em,topsep=3pt,itemsep=2pt,parsep=0pt,label=$\bullet$]\brrulelist
    \item BR-FORUM-38: Only the owner of a comment or an Admin can delete it.
    \item BR-FORUM-39: Comment deletion shall be implemented as a soft delete using the
          \texttt{deletedAt} field.
    \item BR-FORUM-40: Soft-deleted comments shall not be displayed in normal comment listings.
\end{itemize}

\sublabel{Normal Flow}
\begin{enumerate}[leftmargin=1.8em,topsep=3pt,itemsep=2pt,parsep=0pt]
    \item The author or the Administrator opens the post detail screen.
    \item The client renders the bin action on the comments the member may delete.
    \item The member selects the action and the confirmation dialog appears.
    \item The member confirms the deletion.
    \item The client sends the delete request with the access token.
    \item The API loads the comment and confirms that it exists.
    \item The API confirms that the comment belongs to the post of the route.
    \item The API confirms that the requester is the author or an Administrator.
    \item The API writes the deletion timestamp.
    \item The API returns the deleted comment.
    \item The client refreshes the thread, where the comment is no longer present.
\end{enumerate}

\sublabel{Abnormal Flow / Exception Handling}
\begin{flowtable}{Case}{Description}
A1 & The comment does not exist or was already deleted. The API returns
\texttt{Error.ForumCommentNotFound} and the client shows the error message.\\ \hline
A2 & A member tries to delete a comment written by somebody else. The API returns
\texttt{Error.ForumCommentNotOwner} and nothing is deleted.\\ \hline
A3 & An Administrator deletes a comment of another member. The operation is allowed and the
comment disappears from the thread.\\ \hline
A4 & The comment belongs to another post. The request is refused with the not-found error, which
prevents deleting through a foreign route.\\ \hline
A5 & The delete fails at database level. The API returns
\texttt{Error.FailedToDeleteForumComment} and the comment remains visible.\\ \hline
A6 & The member cancels the confirmation. No request is sent and the comment is unchanged.\\ \hline
A7 & The session expires before the confirmation. The single refresh attempt decides between a
successful retry and a redirect to the sign-in page.\\ \hline
\end{flowtable}

\sublabel{Output Data}
\begin{specdatatable}{Data Name}{Description}
\texttt{id} & Identifier of the deleted comment.\\ \hline
\texttt{postId} & Post the comment belonged to.\\ \hline
\texttt{content} & Text of the deleted comment, returned for confirmation.\\ \hline
\texttt{deletedAt} & Server timestamp recorded in the row.\\ \hline
Updated counter & The comment counter of the post, recomputed without the deleted comment.\\ \hline
\end{specdatatable}

\sublabel{Post-Conditions}
\begin{itemize}[leftmargin=1.6em,topsep=3pt,itemsep=2pt,parsep=0pt,label=$\bullet$]
    \item The comment is excluded from the thread and from every comment count.
    \item The row is retained for moderation and for the trash screen of the Administrator.
    \item The post remains published with its remaining comments.
\end{itemize}

% =================================================================

\subsubsection{Function: Like a Comment}

\fnhead{Function Trigger}
A signed-in member selects the heart action on a comment of the thread. A guest is routed to the
sign-in page, because the reaction requires an account.\newline
\textbf{Navigation path:} Home $\rightarrow$ Forum $\rightarrow$ \textit{Category} $\rightarrow$
\textit{Post} $\rightarrow$ Comments $\rightarrow$ Like.

\fnhead{Function Description}

\sublabel{Actors / Roles}
\begin{itemize}[leftmargin=1.6em,topsep=3pt,itemsep=2pt,parsep=0pt,label=$\bullet$]
    \item \textbf{Primary Actor:} Freelancer and Client. An Administrator can also react while
          browsing.
    \item \textbf{Supporting Systems:} Frevia Web application, Frevia BFF proxy, Frevia REST API
          and PostgreSQL through Prisma ORM.
\end{itemize}

\sublabel{Purpose}
\noindent The function allows members to acknowledge a useful reply without writing another
comment. It is a toggle, exactly like the post reaction: the first activation likes the comment
and the second one removes the reaction, so one member can never like the same comment twice.

\sublabel{Interface}
\begin{itemize}[leftmargin=1.6em,topsep=3pt,itemsep=2pt,parsep=0pt,label=$\bullet$]
    \item Heart action in the footer of each comment, followed by the like counter of the
          comment.
    \item Outlined heart when the member has not liked the comment, filled heart when the member
          liked it.
    \item The action is disabled while the request runs.
\end{itemize}

\sublabel{Data Processing Overview}
\begin{itemize}[leftmargin=1.6em,topsep=3pt,itemsep=2pt,parsep=0pt,label=$\bullet$]
    \item The client sends \texttt{POST /api/forums/posts/\{postId\}/comments/\{commentId\}/like}.
    \item The API loads the comment with \texttt{deletedAt} null and rejects the request when it is
          missing.
    \item The API verifies that the comment belongs to the post of the route.
    \item The API looks for a reaction of the member on the comment.
    \item When the reaction exists, the API removes it and answers with \texttt{liked: false}.
    \item When the reaction does not exist, the API inserts it and answers with
          \texttt{liked: true}.
    \item The client invalidates the comment queries of the post, which refreshes the counter and
          the state of the heart of that comment only.
    \item The comment list computes the liked state per comment by matching the reactions of the
          signed-in member, so the states of the other comments are unaffected.
\end{itemize}

\fnhead{Screen Layout (Mock-up Description)}
\screenfigurepair{image/screenshot/sprint1/kvy/BeforeLikeComment.png}{image/screenshot/sprint1/kvy/AfterLikeComment.png}{Comment before and after being liked}
\begin{itemize}[leftmargin=1.6em,topsep=3pt,itemsep=2pt,parsep=0pt,label=$\bullet$]
    \item Before the reaction the heart is outlined and the counter shows the current number of
          likes of the comment.
    \item After the reaction the heart is filled with the accent colour and the counter is
          incremented by one.
    \item Only the selected comment changes; the reactions of the other comments and the likes of
          the post remain untouched.
\end{itemize}

\fnhead{Function Details}

\sublabel{Input Data}
\begin{specdatatable}{Data Name}{Description}
\texttt{postId} & Identifier of the post from the route; it scopes the comment.\\ \hline
\texttt{commentId} & Identifier of the comment to react to.\\ \hline
Access token & Identifies the member who reacts.\\ \hline
Reaction type & Constant value \texttt{like}.\\ \hline
\end{specdatatable}

\sublabel{Validation Rules}
\begin{itemize}[leftmargin=1.6em,topsep=3pt,itemsep=2pt,parsep=0pt,label=$\bullet$]
    \item The requester must be authenticated.
    \item The comment must exist, must not be soft deleted and must belong to the post of the
          route.
    \item A member can hold only one reaction per comment, enforced by the unique constraint of the
          like table on the pair member and comment.
    \item The reaction is attached to the member of the token, never to an identifier sent by the
          client.
\end{itemize}

\sublabel{Business Rules}
\begin{itemize}[leftmargin=1.6em,topsep=3pt,itemsep=2pt,parsep=0pt,label=$\bullet$]\brrulelist
    \item BR-FORUM-41: Only authenticated users can like comments.
    \item BR-FORUM-42: A user can have only one active like reaction per comment.
    \item BR-FORUM-43: The reaction shall be stored in the \texttt{ForumLike} table.
    \item BR-FORUM-44 to BR-FORUM-46: removing a reaction is allowed only when the record exists,
          and it deletes that record.
\end{itemize}

\sublabel{Normal Flow}
\begin{enumerate}[leftmargin=1.8em,topsep=3pt,itemsep=2pt,parsep=0pt]
    \item The signed-in member finds a useful comment in the thread.
    \item The client renders an outlined heart with the like counter of that comment.
    \item The member selects the heart and the client sends the toggle request.
    \item The API confirms that the comment exists and belongs to the post.
    \item The API finds no reaction of this member on this comment.
    \item The API inserts the reaction row.
    \item The API answers with \texttt{liked: true}.
    \item The client refreshes the comment list, fills the heart and increments the counter.
    \item Selecting the heart again removes the reaction and restores the previous counter.
\end{enumerate}

\sublabel{Abnormal Flow / Exception Handling}
\begin{flowtable}{Case}{Description}
A1 & The comment does not exist or was deleted. The API returns
\texttt{Error.ForumCommentNotFound} and the counter is unchanged.\\ \hline
A2 & The comment belongs to another post. The request is refused with the same not-found error.\\ \hline
A3 & The member is not signed in. The client routes to the sign-in page and no request is sent.\\ \hline
A4 & The reaction already exists. The request removes it, which is the intended toggle
behaviour.\\ \hline
A5 & The insert fails at database level. The API returns
\texttt{Error.FailedToLikeForumComment} and the interface keeps the previous state.\\ \hline
A6 & The member clicks the heart twice quickly. The action is disabled while the request runs, so
only one reaction is created.\\ \hline
\end{flowtable}

\sublabel{Output Data}
\begin{specdatatable}{Data Name}{Description}
\texttt{liked} & \texttt{true} when the reaction was added, \texttt{false} when it was removed.\\ \hline
Like row & Member, post, comment, reaction type and creation date of the stored reaction.\\ \hline
Updated counter & The like counter of the comment, recomputed by the client.\\ \hline
\end{specdatatable}

\sublabel{Post-Conditions}
\begin{itemize}[leftmargin=1.6em,topsep=3pt,itemsep=2pt,parsep=0pt,label=$\bullet$]
    \item The comment holds exactly one reaction of the member, or none after the second
          activation.
    \item The counter of the comment matches the number of stored reactions on that comment.
    \item The reactions of the post and of the other comments are unchanged.
\end{itemize}

% =================================================================

\subsubsection{Function: Report Content}

\fnhead{Function Trigger}
A signed-in member who is not the author selects the ``Report'' action under a post or next to a
comment, writes a reason and submits the report. The action is hidden from the author of the
content, because a member cannot report their own post or comment.\newline
\textbf{Navigation path:} Home $\rightarrow$ Forum $\rightarrow$ \textit{Category} $\rightarrow$
\textit{Post} $\rightarrow$ Report.

\fnhead{Function Description}

\sublabel{Actors / Roles}
\begin{itemize}[leftmargin=1.6em,topsep=3pt,itemsep=2pt,parsep=0pt,label=$\bullet$]
    \item \textbf{Primary Actor:} Freelancer and Client.
    \item \textbf{Supporting Systems:} Frevia Web application, Frevia BFF proxy, Frevia REST API
          and PostgreSQL through Prisma ORM; the report is later handled by an Administrator in
          the moderation screen.
\end{itemize}

\sublabel{Purpose}
\noindent The function gives the community a way to flag content that violates the rules, either
a discussion or a single reply. A report never changes the visibility of the content by itself:
it creates a pending moderation entry that the Administrator reviews through the moderation
module.

\sublabel{Interface}
\begin{itemize}[leftmargin=1.6em,topsep=3pt,itemsep=2pt,parsep=0pt,label=$\bullet$]
    \item ``Report'' action with a flag icon under the post body, and the same action in the
          header of each comment.
    \item Dialog titled ``Report Content'' with the prompt ``Please provide a reason for reporting
          this content.''.
    \item ``Reason'' text area with the placeholder ``Why are you reporting this content?''.
    \item ``Cancel'' action and destructive ``Report'' action, which stays disabled while the
          reason is empty.
    \item ``Reported'' badge with a check icon replacing the action once the content has been
          reported by the member.
\end{itemize}

\sublabel{Data Processing Overview}
\begin{itemize}[leftmargin=1.6em,topsep=3pt,itemsep=2pt,parsep=0pt,label=$\bullet$]
    \item When the dialog is opened, the client checks whether the member already reported the
          content through \texttt{GET /api/forums/posts/\{postId\}/is-reported} or
          \texttt{GET /api/forums/posts/\{postId\}/comments/\{commentId\}/is-reported}.
    \item The client sends \texttt{POST /api/forums/posts/\{postId\}/reports} for a post, or
          \texttt{POST /api/forums/posts/\{postId\}/comments/\{commentId\}/reports} for a comment,
          with the reason.
    \item The API verifies that the reported target exists and is not soft deleted.
    \item For a comment, the API also verifies that the comment belongs to the post of the route.
    \item The API looks for an existing report of the same member on the same target and refuses a
          duplicate.
    \item The API inserts the report with the reporter of the token, the target and the reason; the
          status defaults to \texttt{PENDING} and no administrator is assigned yet.
    \item The client shows the confirmation toast, marks the report as done in its cache and closes
          the dialog after a short delay.
    \item From then on the action is replaced by the ``Reported'' badge, and reopening the dialog
          explains that the content is already under review.
\end{itemize}

\fnhead{Screen Layout (Mock-up Description)}
\screenfigure[0.8]{image/screenshot/sprint1/kvy/ReportContent.png}{Report dialog with the reason field}
\begin{itemize}[leftmargin=1.6em,topsep=3pt,itemsep=2pt,parsep=0pt,label=$\bullet$]
    \item The action is available on the post and on every comment, so a member can flag the
          exact item that is problematic.
    \item The dialog explains what is expected and gives a free-text reason, which the moderator
          reads before making a decision.
    \item The confirm action is destructive styled and disabled without a reason, which prevents
          empty reports.
    \item After a successful submission the toast ``Report submitted successfully'' confirms the
          action and the badge ``Reported'' prevents a second submission.
    \item When the content was already reported, the dialog displays a confirmation panel with
          ``You have already reported this content. Our team will review it.''.
\end{itemize}

\fnhead{Function Details}

\sublabel{Input Data}
\begin{specdatatable}{Data Name}{Description}
\texttt{postId} & Identifier of the reported post, or of the post that holds the reported
comment.\\ \hline
\texttt{commentId} & Optional identifier of the reported comment; absent when the post itself is
reported.\\ \hline
\texttt{reason} & Free-text justification of the report, required.\\ \hline
Access token & Identifies the reporter; it becomes \texttt{reporterId} in the stored report.\\ \hline
\end{specdatatable}

\sublabel{Validation Rules}
\begin{itemize}[leftmargin=1.6em,topsep=3pt,itemsep=2pt,parsep=0pt,label=$\bullet$]
    \item The requester must be authenticated.
    \item The reason is required and must contain at least one character.
    \item The reported post must exist and must not be soft deleted, otherwise the request fails
          with \texttt{Error.ForumPostNotFound}.
    \item The reported comment must exist, must not be soft deleted and must belong to the post of
          the route, otherwise the request fails with \texttt{Error.ForumCommentNotFound}.
    \item Exactly one member may report a given target once; a second attempt fails with
          \texttt{Error.ForumReportAlreadyExists}.
    \item The status of a new report is \texttt{PENDING} and the administrator field stays empty
          until a moderator handles it.
\end{itemize}

\sublabel{Business Rules}
\begin{itemize}[leftmargin=1.6em,topsep=3pt,itemsep=2pt,parsep=0pt,label=$\bullet$]\brrulelist
    \item BR-FORUM-47: Only authenticated users can submit reports.
    \item BR-FORUM-48: Reports must be associated with either a post or a comment.
    \item BR-FORUM-49: New reports shall be created with status PENDING.
    \item BR-FORUM-50: A user may submit only one report for the same content item.
\end{itemize}

\sublabel{Normal Flow}
\begin{enumerate}[leftmargin=1.8em,topsep=3pt,itemsep=2pt,parsep=0pt]
    \item The signed-in member reads a post or a comment that violates the community rules.
    \item The member selects the ``Report'' action of that item.
    \item The client opens the dialog and checks whether the item was already reported.
    \item The member writes a reason and confirms.
    \item The client sends the report request with the access token.
    \item The API validates the reason and loads the reported target.
    \item The API confirms that the target exists and, for a comment, that it belongs to the
          post.
    \item The API confirms that the member has not reported the same target yet.
    \item The API inserts the report with status \texttt{PENDING}.
    \item The API returns the created report.
    \item The client shows the success toast, marks the cache as reported and closes the dialog.
    \item The ``Reported'' badge replaces the action for this member.
\end{enumerate}

\sublabel{Abnormal Flow / Exception Handling}
\begin{flowtable}{Case}{Description}
A1 & The reason is empty. The confirm action is disabled, and a forced request is rejected by the
validation pipe.\\ \hline
A2 & The post was deleted before the submission. The API returns
\texttt{Error.ForumPostNotFound} and no report is stored.\\ \hline
A3 & The comment was deleted before the submission, or does not belong to the post. The API returns
\texttt{Error.ForumCommentNotFound}.\\ \hline
A4 & The member already reported the same item. The API returns
\texttt{Error.ForumReportAlreadyExists} and the client displays the message and switches to the
reported state.\\ \hline
A5 & An Administrator opens the moderation screen. Listing and updating reports require the
Administrator role, otherwise the API returns \texttt{Error.ForumReportForbidden}.\\ \hline
A6 & The insert fails at database level. The API returns
\texttt{Error.FailedToCreateForumReport} and the dialog stays open with the reason.\\ \hline
A7 & The member is not signed in. The action is not rendered for guests, and a direct request is
rejected with an authentication error.\\ \hline
A8 & The member reports their own content through a crafted request. The record can be created,
but the action is never offered in the interface, and the moderation screen treats it like any
other report.\\ \hline
\end{flowtable}

\sublabel{Output Data}
\begin{specdatatable}{Data Name}{Description}
\texttt{id} & Identifier of the created report.\\ \hline
\texttt{reporterId} & Member who submitted the report.\\ \hline
\texttt{postId}, \texttt{commentId} & Reported targets; exactly one of them identifies the item,
the other may be null for a comment report.\\ \hline
\texttt{reason} & Reason written by the reporter.\\ \hline
\texttt{status} & \texttt{PENDING} for a new report; later \texttt{REVIEWED},
\texttt{RESOLVED} or \texttt{DISMISSED} after the moderator decision.\\ \hline
\texttt{adminId}, \texttt{createdAt} & Administrator who handled the report, null at creation, and
the creation timestamp.\\ \hline
\end{specdatatable}

\sublabel{Post-Conditions}
\begin{itemize}[leftmargin=1.6em,topsep=3pt,itemsep=2pt,parsep=0pt,label=$\bullet$]
    \item One pending report exists for the reported item and the reporter.
    \item The reported content is not modified and stays visible until a moderator decides.
    \item The report appears in the moderation queue of the Administrator, where it can be reviewed,
          resolved or dismissed.
    \item The reporter sees the ``Reported'' state and cannot submit the same report again.
\end{itemize}
